\documentclass[11pt]{article}
\usepackage[T1]{fontenc}
\usepackage{microtype}
\usepackage[a4paper,margin=29mm]{geometry}
\usepackage{amsmath,amssymb,amsthm,mathtools}
\usepackage[colorlinks=true,linkcolor=blue,citecolor=blue,urlcolor=blue]{hyperref}
\hypersetup{pdftitle={Cubic sphere moments and the two-pair Image problem},pdfauthor={Roy van Rijn}}
\newcommand{\C}{\mathbb C}
\newcommand{\A}{\mathbb A}
\newcommand{\E}{\mathcal E_2}
\newcommand{\CT}{\operatorname{CT}}
\newcommand{\dd}{\,d}
\DeclareMathOperator{\Res}{Res}
\newtheorem{theorem}{Theorem}[section]
\newtheorem{proposition}[theorem]{Proposition}
\newtheorem{lemma}[theorem]{Lemma}
\newtheorem{corollary}[theorem]{Corollary}
\theoremstyle{remark}\newtheorem{remark}[theorem]{Remark}
\title{Cubic sphere moments and the two-pair Image problem}
\author{Roy van Rijn\\\small ORCID: \href{https://orcid.org/0009-0000-2808-0645}{0009-0000-2808-0645}}
\date{}
\begin{document}
\maketitle
\begin{abstract}
We classify the polynomials of degree at most three on the complex affine
unit sphere whose positive spherical moments all vanish. Up to an
orthogonal change of coordinates and a nonzero scalar, they have strictly
one-sided circular support or equal \(y+xt-x^3/2\), where
\(t^2+2xy=1\). Their mixed moments vanish for exponents greater than three
times the degree of the multiplier, and this coefficient is optimal.
The proof combines a proper-map obstruction for exponential volume
periods, explicit resolutions at infinity, and a small exact moment
elimination. It implies the corresponding sharp bound for every
bihomogeneous polynomial of bidegree \((3,3)\) in two contraction pairs.
An explicit bidegree-\((4,4)\) counterexample shows that four is the
minimum failing balanced bidegree.
\end{abstract}

\section{The result and its mechanism}

Let \(A=\C[x,y,t]/(t^2+2xy-1)\), with the filtration \(A_{\leq d}\)
induced by ordinary polynomial degree. Write \(q=t^2+2xy\).
The normalized spherical integral is the linear functional \(L:A\to\C\)
characterized by
\begin{equation}\label{eq:moments}
 L(x^a y^b t^c)=
 \begin{cases}
 \displaystyle\frac{a!(2k-1)!!}{(2a+2k+1)!!},&a=b,\ c=2k,\\[4pt]
 0,&\text{otherwise},
 \end{cases}\qquad (-1)!!=1.
\end{equation}
Indeed, put \(x=(u_1+iu_2)/\sqrt2\),
\(y=(u_1-iu_2)/\sqrt2\), and \(t=u_3\), and average over the real
unit sphere in \(\mathbb R^3\). The functional is invariant under
\(O(q,\C)\). We call \(p\) \emph{pure} if \(L(p^m)=0\) for every
integer \(m\geq1\).
Under the circular action
\((x,y,t)\mapsto(\lambda x,\lambda^{-1}y,t)\), the weights of
\(x,y,t\) are \(1,-1,0\). A polynomial has strictly positive support
if every nonzero monomial in its normal form has positive weight;
the zero polynomial is included.

\begin{theorem}\label{thm:sphere}
A polynomial \(p\in A_{\leq3}\) is pure if and only if, after an
orthogonal change of coordinates, it has strictly positive support or is
\[
 \alpha\left(y+xt-\frac{x^3}{2}\right),\qquad \alpha\in\C^*.
\]
For every \(g\in A_{\leq D}\), purity implies
\begin{equation}\label{eq:cutoff}
 L(gp^m)=0\qquad(m>3D).
\end{equation}
The coefficient three is optimal: for every \(D\geq1\),
\begin{equation}\label{eq:sharp}
 L\left(y^D\left(y+xt-\frac{x^3}{2}\right)^{3D}\right)
 =\frac{(-1)^D2^{2D}((3D)!)^2}{(6D+1)!}\ne0.
\end{equation}
\end{theorem}

The one-sided case is immediate from circular weights. The second case
is also geometric, but uses a nonlinear change of coordinates. The map
\begin{equation}\label{eq:shear}
 \phi(x)=x,\qquad \phi(t)=t-x^2,\qquad
 \phi(y)=y+xt-\frac{x^3}{2}
\end{equation}
is an automorphism of the sphere and preserves \(L\).
Thus all powers of \(\phi(y)\) integrate to zero. Its inverse has degree
three, explaining the coefficient in \eqref{eq:cutoff}. The difficult
point is proving that these exhaust the cubics.

Here is the reason for that exhaustion. Pure moments make the
exponential period \(\int e^{-\tau p}\Omega\) identically one. If the
sphere volume extends to a smooth compactification on which \(p\) is a
proper map to the affine line, a positive \(V\)-filtration obstruction
rules out that constant period. We resolve the graph at infinity and
identify exactly the boundary configurations that prevent this
extension. They force a coordinate frame in which \(p=ay+G(x,t)\),
and, when \(a\ne0\), the cubic term \(t^3\) can be excluded. Eight
explicit moments then give the remaining normal form. The geometric
step is essential: a finite moment calculation in one coordinate chart
does not classify arbitrary cubics.

For the Image problem, let \(k\) be a field of characteristic zero and
set \(B=k[\xi_1,\xi_2,z_1,z_2]\). Define
\[
 \E(\xi^\alpha z^\beta)=\partial_z^\alpha z^\beta,
 \qquad
 \mathcal M_2=\sum_{j=1}^2(\partial_{z_j}-\xi_j)B.
\]
The identity \(\mathcal M_2=\ker\E\) is the standard contraction
formulation of the Special Image problem
\cite[Section~3]{EWZ11}; a direct proof is given below.
A polynomial of bidegree \((d,d)\) is homogeneous of degree \(d\)
separately in \(\xi\) and \(z\).

\begin{theorem}\label{thm:image}
Let \(F\in B\) have bidegree \((3,3)\), and suppose
\(\E(F^m)=0\) for every \(m\geq1\). Then, for every nonzero \(Q\in B\),
\[
 \E(QF^m)=0\qquad(m>3\deg_z Q).
\]
The coefficient three is optimal. The minimum balanced bidegree
\(d\) admitting a counterexample to eventual mixed-contraction
vanishing is four.
\end{theorem}

The Image framework of van den Essen, Wright and Zhao
\cite{EWZ11} concerns arbitrary polynomials and does not supply the
cubic classification above. Long's counterexamples to the mixed
interval--torus moment conjecture and the Mathieu conjecture for
\(SU(2)\) \cite{Long26} establish failures outside restricted
classes. Our positive result concerns the full filtered space
\(A_{\leq3}\), or, equivalently, fixed balanced bidegree three in two
pairs. Long's phase-balanced Hopf witness \cite[Theorem~3.1]{LongHopf26} gives the quartic upper boundary. We reproduce its coefficient proof in our normalization to make the degree
boundary explicit. We make no assertion for arbitrary sums of different
bidegrees or for more contraction pairs.

\section{Spherical integration and contraction}\label{sec:bridge}

On the circular chart \(y=(1-t^2)/(2x)\), spherical integration becomes
\begin{equation}\label{eq:ct}
 L(a)=\frac12\int_{-1}^1
 \CT_x\,a\left(x,\frac{1-t^2}{2x},t\right)\dd t.
\end{equation}
The beta integral proves \eqref{eq:moments}. The relation \(q=1\)
gives a unique normal form with no monomial divisible by \(xy\).
Its monomials of degree at most \(d\) number \((d+1)^2\).
Orthogonal invariance follows either from the Cartesian sphere
moments or from their contraction formula: odd tensors vanish and
even tensors are sums of pairings of \(q\).

Introduce the bilinears
\begin{equation}\label{eq:bilinears}
 R=\xi_1z_1+\xi_2z_2,\quad X=\xi_1z_2,\quad
 Y=2\xi_2z_1,\quad T=\xi_1z_1-\xi_2z_2.
\end{equation}
They satisfy \(T^2+2XY=R^2\). A balanced polynomial of bidegree
\((n,n)\) is a homogeneous degree-\(n\) polynomial in these bilinears,
modulo that relation. To see surjectivity, pair the \(n\) dual factors
with the \(n\) coordinate factors in each monomial. Both spaces have
dimension \((n+1)^2\), so restriction to \(R=1\) identifies this
space with \(A_{\leq n}\).

\begin{lemma}\label{lem:bridge}
If \(H\) has bidegree \((n,n)\) and restricts to \(h\in A_{\leq n}\),
then
\begin{equation}\label{eq:bridge}
 \E(H)=(n+1)!L(h).
\end{equation}
Moreover \(\mathcal M_2=\ker\E\).
\end{lemma}
\begin{proof}
For the first identity, substitute \(\xi_j=\overline z_j\) on the
unit sphere in \(\C^2\). Its normalized integral sends
\(\overline z^\alpha z^\beta\), with \(|\alpha|=|\beta|=n\), to
\(\delta_{\alpha\beta}\alpha!/(n+1)!\). This follows by phase
integration and \(\int_0^1s^{\alpha_1}(1-s)^{\alpha_2}\dd s
=\alpha!/(n+1)!\). The resulting variables \((X,Y,T)\) have the
normalized sphere moments \eqref{eq:moments}: their circular phase is
uniform and \(T=2|z_1|^2-1\) is uniform on \([-1,1]\).

Contraction kills each \((\partial_{z_j}-\xi_j)H\), so
\(\mathcal M_2\subset\ker\E\). Conversely the relation
\(\xi_j H\equiv\partial_{z_j}H\pmod{\mathcal M_2}\), iterated
on monomials, gives \(H\equiv\E(H)\pmod{\mathcal M_2}\).
Thus \(\ker\E\subset\mathcal M_2\).
\end{proof}

In Cartesian coordinates the normalized algebraic volume is
\[
 \Omega=\frac1{2\pi}\Res_{q=1}
 \frac{du_1\wedge du_2\wedge du_3}{q-1}.
\]
Its restriction to the oriented real sphere \(\Sigma\) is normalized
area, and \(L(a)=\int_\Sigma a\Omega\). The entire function
\(\int_\Sigma e^{-\tau p}\Omega\) is identically one precisely when
\(p\) is pure.

\section{A proper volume obstruction}\label{sec:proper-volume}

We use one consequence of the Kashiwara--Malgrange filtration. It applies
to critical curves as well as isolated critical points.

\begin{lemma}\label{lem:proper-volume}
Let \(f:Y\to\mathbb A^1\) be a proper nonconstant morphism from a smooth
irreducible complex algebraic variety, and let \(\omega\) be a global
regular top form on \(Y\). If \(\gamma\) is a fixed compact cycle of real
dimension \(\dim_{\mathbb C}Y\), then
\[
 \int_\gamma e^{-\tau f}\omega
\]
is not a nonzero constant.
\end{lemma}

\begin{proof}
Write \(t\) for the coordinate on the target, factor \(f\) as its graph
\(i:Y\hookrightarrow Y\times\mathbb A^1_t\) followed by the projection
\(\pi\), and put \(\mathcal B_f=i_+\mathcal O_Y\).
We use the decreasing filtration \(V^\alpha\), normalized by the
nilpotence of \(\partial_t t-\alpha\) on \(\operatorname{gr}_V^\alpha\).
Negativity of the local Bernstein--Sato roots
\cite[Corollary~5.2]{Kas76} gives
\begin{equation}\label{eq:geo-star}
\delta(t-f)\in V^{>0}\mathcal B_f
 \qquad\text{\cite[Corollary~5.15 and Example~5.19]{CDM24}}.
\end{equation}
Nonconstancy ensures that every local germ of \(f\) is nonzero; no
isolated-critical-point hypothesis is used.

Let \(M=\mathcal H^0f_+\mathcal O_Y\), and let \(m\) be the image of
\(\omega\). This is a global top-form class in the relative de Rham
\emph{hypercohomology} computing \(M\), not a quotient of global forms
by global exact forms. Near the compact fiber \(f^{-1}(0)\), finitely
many local bounds in \eqref{eq:geo-star} give a common \(\alpha>0\) for
\(\omega\otimes\delta(t-f)\). Properness allows the target neighborhood
to be shrunk so that this finite cover suffices. The projection \(\pi\)
is proper on \(\operatorname{Supp}\mathcal B_f\).
Compatibility of \(V^\alpha\) with this derived direct image
\cite[Theorem~8.1]{CDM24} therefore gives
\begin{equation}\label{eq:geo-starstar}
m\in V^\alpha M\quad\text{for some }\alpha>0.
\end{equation}
This tracks the given form itself, with no parameter-dependent change
of seed.

Take global sections on the affine target and Fourier transform by
\(t\mapsto-\partial_\tau\), \(\partial_t\mapsto\tau\).
Compact-cycle integration is the de Rham period pairing: Stokes
annihilates relative exact classes, and differentiation in \(\tau\)
inserts \(-f\). Hence it respects every Weyl-algebra relation of the
Fourier image \(\widehat m\).
If its period were a nonzero constant, after scaling the pairing the
cyclic submodule
\[
 \widehat N=\mathbb C[\tau]\langle\partial_\tau\rangle\widehat m
 \subset\widehat M
\]
would have a quotient to \(\mathcal O_{\mathbb A^1_\tau}\) sending
\(\widehat m\) to \(1\).
Inverse Fourier transform gives a submodule \(N\subset M\) and a
quotient \(N\to\delta_0=D_{\mathbb A^1_t}/D_{\mathbb A^1_t}t\)
sending \(m\) to its nonzero generator.
But submodules carry the induced \(V\)-filtration and morphisms
preserve it, whereas \(V^\alpha\delta_0=0\) for \(\alpha>0\)
\cite[Proposition~3.19(ii), Corollary~3.20, Example~3.18]{CDM24}.
This contradicts \eqref{eq:geo-starstar}. The cyclic quotient was constructed before
localization, so no assertion about a localized seed is needed.
\end{proof}

\section{Reduction to a constant profile}\label{sec:constant-profile}

Put \(X=\{t^2+2xy=1\}\) and \(A=\mathbb C[x,y,t]/(t^2+2xy-1)\).
The normalized residue volume \(\Omega\) on \(X\) has
\(\int_\Sigma\Omega=1\) on the real unit sphere in Cartesian
coordinates. Thus \(L(a)=\int_\Sigma a\Omega\).
In a circular chart it is a nonzero constant multiple of
\(dy\wedge dt/(2y)\). Constant factors will be suppressed in the
local regularity calculations below.

\begin{proposition}\label{prop:constant-profile}
If \(p\in A_{\leq3}\) satisfies \(L(p^m)=0\) for every \(m\geq1\),
then in some orthogonal coordinate frame
\[
 p=ay+G(x,t),\qquad \deg G\leq3,\qquad a\,[t^3]G=0.
\]
\end{proposition}

\begin{proof}
Choose a representative \(P=P_d+\cdots+P_0\) of minimal degree \(d\),
with homogeneous pieces \(P_j\). The cases \(d\leq1\) already have
the stated form. For \(d\geq2\), minimality implies
\(P_d|_{\{q=0\}}\neq0\): otherwise \(P_d=qH\), and subtracting
\((q-1)H\) lowers the degree.

Compactify \(X\) by
\[
 \bar X=\{T^2+2XY=R^2\}\subset\mathbb P^3.
\]
The actual closure \(\Gamma_d\) of the graph of \(p\) is proper over
the phase line and is given by
\begin{equation}\label{eq:geo-1}
\sum_{j=0}^d R^{d-j}P_j(X,Y,T)-sR^d=0.
\end{equation}
There is no component supported on \(R=0\), because the nonzero
restriction of \(P_d\) has only finitely many zeros on the boundary
conic. All modifications below are strict transforms of this
closure, and are isomorphisms over \(X\).

Move a boundary point to \([X:Y:T:R]=[0:1:0:0]\) by an orthogonal
change of frame, and use coordinates
\begin{equation}\label{eq:geo-2}
(X,Y,T,R)=((z^2-u^2)/2,1,u,z).
\end{equation}
The original volume is, up to a nonzero constant,
\begin{equation}\label{eq:geo-3}
\Omega=\frac{du\wedge dz}{2z^2}.
\end{equation}

Suppose first that \(d=2\).
At a boundary root, the coefficient of \(y^2\) in \(P_2\) is zero.
If the root of \(P_2|_{\{q=0\}}\) is repeated, its derivative in
\(u\) also vanishes, so the coefficient of \(yt\) is zero.
Replacing \(xy\) by \((1-t^2)/2\) then gives
\(p=ay+G(x,t)\), as required.
If no such frame exists, all boundary roots are simple.
The local equation \(F=0\) in \eqref{eq:geo-1} is then smooth at every boundary
point, with \(F_s=-z^2\). Its regular dualizing residue \(\eta_F\)
satisfies
\[
 \Omega=-\tfrac12\eta_F.
\]
Consequently \(\Omega\) extends to the whole smooth proper graph.
Lemma~\ref{lem:proper-volume} contradicts
\(\int_\Sigma e^{-\tau p}\Omega=1\).
This settles \(d=2\), without a separate quadratic classification.

Now let \(d=3\), and suppose that no frame satisfies the conclusion.
At a boundary point in \eqref{eq:geo-2}, its equation has derivatives
\[
 F_u(0,0,s)=[y^2t]P_3,\qquad
 F_z(0,0,s)=[y^2]P_2,
\]
while the root condition is \([y^3]P_3=0\).
If either derivative is nonzero, the graph is smooth there.
Since \(F_s=-z^3\), its dualizing residue satisfies
\begin{equation}\label{eq:geo-4}
\Omega=-\tfrac z2\,\eta_F,
\end{equation}
which extends regularly across that point.

It remains to treat points where both derivatives vanish.
If \(C=[xy^2]P_3\), replace \(P\) by
\(P-\frac C2y(q-1)\). This removes \(xy^2\) and changes neither
\(p\) nor the graph on \(\bar X\). The representative is now affine
linear in \(y\). Eliminating its remaining \(xy\) factors gives
\[
 p=y(a+bt+ct^2)+G(x,t),\qquad \deg G\leq3.
\]
In \eqref{eq:geo-2} the graph equation is
\begin{equation}\label{eq:geo-5}
F=cu^2+buz+az^2+R_{\geq3}(u,z)-sz^3=0.
\end{equation}
Indeed \(x^it^j\) in \(G\) contributes
\(2^{-i}(z^2-u^2)^iu^jz^{3-i-j}\), of total order at least \(3+i\).
A constant profile with \(a[t^3]G=0\) would finish the proof. We resolve the other cases.

If \(b^2-4ac\neq0\), blow up the boundary section.
In the chart \(u=zv\), the strict transform is \(F=z^2F_1\), with
\[
 F_1=cv^2+bv+a+zK(v,z,s),\qquad (F_1)_s=-z.
\]
Its exceptional intersections are the simple roots of the profile,
so it is smooth there, and
\begin{equation}\label{eq:geo-6}
\Omega=\frac{dv\wedge dz}{2z}=-\tfrac12\eta_{F_1}.
\end{equation}
If \(c=0,b\neq0\), the additional tangent direction is covered by
\(z=uv\). Here \(F=u^2F_2\), and the exceptional point \(u=v=0\)
has \((F_2)_v=b\neq0\). Moreover
\begin{equation}\label{eq:geo-7}
(F_2)_s=-uv^3,\qquad
 \Omega=\frac{du\wedge dv}{2uv^2}=-\tfrac v2\,\eta_{F_2}.
\end{equation}
Thus both charts are smooth along their full exceptional boundary
and the volume is regular there.

The remaining nonconstant profile has \(c\neq0\) and
\(b^2-4ac=0\). Scale \(c=1\), write the profile as \((t-r)^2\),
and put \(w=u-rz\). Blow up in the chart \(w=zv\), which contains
the entire exceptional intersection. The strict equation is
\begin{equation}\label{eq:geo-8}
F_1=v^2+z(k(v)-s)+z^2\rho(v,z)=0.
\end{equation}
It is smooth on the boundary except possibly at
\((v,z,s)=(0,0,k(0))\).
The Hessian determinant at that point is \(-2\), independently of
\(k\) and \(\rho\); it is an ordinary \(A_1\) surface double point.
By \eqref{eq:geo-6}, \(\Omega\) is a constant multiple of its dualizing generator.
Blow up this point. The strict transform is smooth, with a smooth
conic exceptional divisor, and the resolution is crepant:
the ambient point blow-up adds twice the exceptional divisor to
the canonical divisor, while the double-point hypersurface has
multiplicity two. Adjunction cancels these contributions.
The given volume therefore extends holomorphically over the
resolution.


For a constant profile, our contrary assumption says \(a\neq0\) and
\(h=[t^3]G\neq0\). Write
\(G(0,t)=ht^3+g_2t^2+g_1t+g_0\).
Equation \eqref{eq:geo-5} is then
\begin{equation}\label{eq:geo-9}
F=az^2+hu^3+g_2u^2z+g_1uz^2+(g_0-s)z^3+R_{\geq4}(u,z).
\end{equation}
Blow up in the chart \(z=uv\). The strict transform is
\begin{equation}\label{eq:geo-10}
F_1=av^2+u\bigl(h+g_2v+g_1v^2+(g_0-s)v^3\bigr)
          +u^2H(u,v)=0.
\end{equation}
The exceptional intersection is the single point \(u=v=0\), where
\((F_1)_u=h\neq0\). Hence it is smooth for every finite \(s\).
Furthermore \((F_1)_s=-uv^3\), and \eqref{eq:geo-7} again gives
\(\Omega=-v\eta_{F_1}/2\), a regular form.
Equivalently, solving \(u=u(v,s)\) gives
\[
 \Omega=\pm\frac{v}{2(F_1)_u}\,dv\wedge ds.
\]
The complementary chart \(u=zv\) has strict equation \(a+O(z)=0\)
and therefore no exceptional intersection. These charts cover the
complete exceptional divisor; no boundary-supported component
has been retained.

There are finitely many boundary points. The indicated blow-ups
of boundary sections, followed by the finitely many ordinary
double-point resolutions, give a smooth irreducible surface \(Y\)
proper over \(\mathbb A^1_s\), containing \(X\), on which the original
volume \(\Omega\) is regular. The phase is nonconstant and the
original sphere remains a compact cycle of nonzero volume.
Purity gives \(\int_\Sigma e^{-\tau p}\Omega=1\), contradicting
Lemma~\ref{lem:proper-volume}.
A frame with constant profile and \(a[t^3]G=0\) must therefore exist.
\end{proof}

\section{The finite moment calculation}\label{sec:finite}

Write \(J(f)=\frac12\int_{-1}^1 f(t)\,dt\).

\begin{lemma}[Interval cubics]
\label{lem:interval-cubic}
If \(u\in\mathbb C[t]\), \(\deg u\le3\), and
\(J(u^m)=0\) for \(1\le m\le4\), then \(u=0\).
\end{lemma}

\begin{proof}
The first equation puts
\(u=h t^3+v(t^2-1/3)+ct\). Write \(M_j=J(u^j)\). Direct integration gives
\begin{align}
 M_2&=\frac{105c^2+126ch+45h^2+28v^2}{315},\notag\\
 M_3\bmod M_2&=\frac{8v(56c^2+105ch+45h^2)}{2205},\notag\\
 M_4\bmod M_2&=-\frac{16R(c,h)}{735735},
 \label{eq:interval-cubic-remainders}
\end{align}
where remainders are taken as polynomials in \(v\), and
\[
 R(c,h)=28028c^4+84084c^3h+93275c^2h^2+46410ch^3+8955h^4.
\]
If \(h=0\), division of \(M_3\) by \(M_2\), now in the variable
\(c\), has remainder \(-256v^3/4725\). Thus \(v=c=0\).

Suppose \(h\ne0\). Scale \(h=1\), rename \(c=z\), and put
\[
 P=56z^2+105z+45,\qquad B=105z^2+126z+45,\qquad R=R(z,1).
\]
If \(v\ne0\), the equations force \(P=R=0\). But
\[
 R\bmod P=-\frac{315}{128}(299z+159),\qquad
 P(-159/299)=\frac{446976}{89401}\ne0.
\]
If \(v=0\), they force \(B=R=0\), whereas
\[
 R\bmod B=-\frac{84}{125}(117z-110),\qquad
 B(110/117)=\frac{1169375}{4563}\ne0.
\]
Both alternatives are impossible.
\end{proof}

The geometric reduction supplies either a phase independent of \(y\),
or a constant profile in which the phase-zero term \(t^3\) is absent.
The first case is already settled by Lemma~\ref{lem:interval-cubic}:
if \(p\in\mathbb C[x,t]\), phase averaging leaves its coefficient of
\(x^0\); purity makes that interval cubic zero. All remaining terms
have strictly positive circular weight.

\begin{proposition}[Constant-profile normal forms]
\label{prop:finite}
Let
\begin{equation}
 p=ay+(v_0+v_1x)t^2+(c+dx+ex^2)t+\sum_{j=0}^3r_jx^j.
 \label{eq:constant-profile-family}
\end{equation}
If \(L(p^m)=0\) for \(1\le m\le8\), then precisely one of the
following cases holds:
\begin{enumerate}
\item \(a=0\) and \(v_0=c=r_0=0\);
\item \(a\ne0\), \(v_0=v_1=e=0\), and
\begin{equation}
 \sum_{j=0}^3r_jx^j=-\frac{x(c+dx)^2}{2a}.
 \label{eq:constant-profile-shear}
\end{equation}
\end{enumerate}
\end{proposition}

\begin{proof}
If \(a=0\), phase averaging leaves \(v_0t^2+ct+r_0\), which vanishes
by Lemma~\ref{lem:interval-cubic}. This is precisely the first case.

If \(a\ne0\), divide by \(a\), rename coefficients, and assume
\(a=1\). The orthogonal null rotation
\begin{equation}
 x=x',\qquad t=t'+kx',\qquad y=y'-kt'-k^2x'/2
 \label{eq:constant-profile-null-rotation}
\end{equation}
preserves \eqref{eq:constant-profile-family}, leaves \(v_0,v_1\)
unchanged, and replaces \(c\) by \(c-k\). Take \(k=c\).

Here is a closed formula for the finite coefficient calculation. For
\[
 p=y+\sum_{\nu=1}^N a_\nu x^{i_\nu}t^{j_\nu},\qquad i_\nu\ge0,
\]
let \(\alpha\in\mathbb N^N\),
\(k=\sum i_\nu\alpha_\nu\), and \(j=\sum j_\nu\alpha_\nu\).
Phase balance forces exactly \(k\) factors of \(y\), so
\begin{equation}
 L(p^m)=
 \sum_{\substack{\sum(1+i_\nu)\alpha_\nu=m\\j\ \mathrm{even}}}
 \frac{m!(j-1)!!}{(2k+j+1)!!\prod_\nu\alpha_\nu!}
 \prod_\nu a_\nu^{\alpha_\nu}.
 \label{eq:constant-profile-coefficients}
\end{equation}
The factorial for the \(y\)-factors cancels the \(k!\) in
\eqref{eq:moments}. Thus this formula computes every
coefficient over \(\mathbb Q\), for all parameter values at once.

For \(a=1,c=0\), the first four equations solve successively for
\(r_0,r_1,r_2,r_3\). Their respective coefficients are the nonzero
constants \(1,2/3,2/5,8/35\), and the solutions are
\begin{align}
 r_0&=-v_0/3,&
 r_1&=-(2v_0^2+3v_1)/15,\notag\\
 r_2&=-4v_0(31v_0^2+54v_1)/945,\notag\\
 r_3&=-\frac{14175d^2+4624v_0^4+10152v_0^2v_1+3024v_1^2}{28350}.
 \label{eq:constant-profile-triangular}
\end{align}
After substitution, if \(v_0=0\), the eighth equation is
\(-65536v_1^4/41023125=0\). Then the sixth equation becomes
\(8e^2/231=0\). Hence \(v_1=e=0\), and
\eqref{eq:constant-profile-triangular} leaves only \(r_3=-d^2/2\).

Suppose \(v_0\ne0\), and set
\[
 b=v_1/v_0^2,\qquad D=d^2/v_0^4,\qquad E=de/v_0^5,\qquad F=e^2/v_0^6.
\]
Let \(N_i\), \(5\le i\le8\), be the primitive integer numerators of
the four residual moments, with positive coefficient of \(v_0^i\).
The equations \(N_i/v_0^i=0\) are linear in \(D,E,F\). Their augmented
matrix, in the column order \(D,E,F,1\), is
\begin{equation}
 M(b)=\begin{pmatrix}
 155925&155925&0&B_5\\
 77395500-116093250b&77395500&-58046625&B_6\\
 226800-297675b&226800+127575b&-212625&B_7\\
 4819500-216877500b&4819500-65063250b&-75907125&B_8
 \end{pmatrix},
 \label{eq:constant-profile-matrix}
\end{equation}
where
\begin{align*}
 B_5&=35504+94992b+52704b^2,\\
 B_6&=30907712+63343584b+3500496b^2-12573792b^3,\\
 B_7&=53504+17328b-211968b^2-108864b^3,\\
 B_8&=7986496-27544704b-90351504b^2-16632864b^3+9144576b^4.
\end{align*}
Its coefficient columns satisfy \(C_D-C_E-2bC_F=0\). Consistency
therefore forces every three-by-three augmented minor to vanish.
Columns \(D,F,1\), with the first three rows and with rows one, two,
four, give, up to nonzero rational factors,
\begin{align*}
 f(b)&=2485161b^3+6647562b^2+4908168b+1097984,\\
 g(b)&=10216206b^4+24356619b^3-25372278b^2-41158800b-11711872.
\end{align*}
These polynomials are coprime over \(\mathbb Q\): their integer
Sylvester determinant is congruent to \(2\) modulo \(19\). A nonzero
integer determinant has nonzero image in \(\mathbb C\), so no common
root exists. This excludes \(v_0\ne0\), without dividing by \(b\)
or assuming any further coefficient is nonzero. Undoing
\eqref{eq:constant-profile-null-rotation} and the scalar normalization
gives \eqref{eq:constant-profile-shear}.

\end{proof}

\paragraph{Exact coefficient verification.}
All finite inputs and identities used above are displayed in
\eqref{eq:interval-cubic-remainders}--\eqref{eq:constant-profile-matrix}
and the accompanying formulas. The verification program accompanying
the article derives the eight moments from
\eqref{eq:constant-profile-coefficients}, reproduces the triangular
solutions and matrix, and evaluates the seven-by-seven Sylvester
determinant over \(\mathbb F_{19}\). It uses exact arithmetic and
requires no search over the coefficient family or external polynomial
data. Zeroing a row of the determinant is rejected. The program is a
replay of the displayed finite argument; the all-order conclusions
follow from the resulting normal forms.


\section{Normal forms and the sharp cutoff}\label{sec:normal-forms}

The geometric reduction and the finite calculation show that a pure
cubic has positive support or has the form
\begin{equation}\label{eq:family}
 p=ay+(c+dx)t-\frac{x(c+dx)^2}{2a},\qquad a\ne0.
\end{equation}
The orthogonal null rotation
\[
 x\mapsto x,\quad t\mapsto t+kx,\quad
 y\mapsto y-kt-\frac{k^2x}{2},\qquad k=c/a,
\]
removes \(c\). If \(d=0\), the result is \(ay\), which becomes
positive by interchanging \(x\) and \(y\). If \(d\ne0\), a nonzero
scalar and the torus \((x,y)\mapsto(\lambda x,\lambda^{-1}y)\)
normalize the coefficients of \(y\) and \(xt\) to one, giving the
exceptional form in Theorem~\ref{thm:sphere}.

For completeness, \eqref{eq:shear} preserves \(q\) and is
\(\phi=\exp(-x\delta)\), where
\(\delta=x\partial_t-t\partial_y\) is locally nilpotent.
Orthogonal invariance gives \(L(\delta f)=0\). Since \(\delta x=0\),
it also gives \(L(x\delta f)=0\), and the finite exponential series
shows \(L(\phi f)=L(f)\). Thus the exceptional form is pure.
If \(g\in A_{\leq D}\), then \(\phi^{-1}g\) has degree at most
\(3D\). Circular weights give
\[
 L(g\phi(y)^m)=L((\phi^{-1}g)y^m)=0\qquad(m>3D).
\]
For positive support, the same weight argument already gives zero for
\(m>D\). This proves the classification and \eqref{eq:cutoff}.

To verify sharpness, apply \(\phi^{-1}\) to \(y^D\phi(y)^{3D}\).
Now \(\phi^{-1}(y)=y-xt-x^3/2\). In
\((y-xt-x^3/2)^D y^{3D}\), the only zero-weight term is
\((-x^3/2)^D y^{3D}\). Formula \eqref{eq:moments} gives
\eqref{eq:sharp}.

\begin{proof}[Proof of Theorem~\ref{thm:image}]
First work over \(\C\), and let \(p\in A_{\leq3}\) correspond to
\(F\). Lemma~\ref{lem:bridge} shows that \(p\) is pure.
If \(Q\) has bidegree \((b,b)\), its restriction has degree at most
\(b\), so \eqref{eq:cutoff} and the same lemma prove the claim.

For general bihomogeneous \(Q\) of bidegree \((a,b)\), the contraction
is automatically zero when \(a>b\). If \(a\leq b\), put
\(r=b-a\). The polynomial \(P_m=\E(QF^m)\) is homogeneous of
\(z\)-degree \(r\). For every multi-index \(|\gamma|=r\),
\[
 \partial_z^\gamma P_m=\E(\xi^\gamma QF^m)=0\qquad(m>3b),
\]
because \(\xi^\gamma Q\) is balanced of bidegree \((b,b)\).
All coefficients of \(P_m\) therefore vanish. Decomposing an arbitrary
\(Q\) into its bidegrees proves the bound as stated.

Every finite collection of characteristic-zero coefficients generates
a field admitting an embedding into \(\C\). Applying the proved result
there and using injectivity proves the assertion over \(k\).
For sharpness use
\[
 F_*=R^2Y+RXT-X^3/2,\qquad Q=Y^D.
\]
The nonzero value at \(m=3D\) is \((10D+1)!\) times
\eqref{eq:sharp}. The same argument for bidegrees \((d,d)\),
\(d\leq3\), follows from \(A_{\leq d}\subset A_{\leq3}\).
The quartic example below completes the degree boundary.
\end{proof}

\section{The quartic boundary}\label{sec:quartic}

The following is Long's Hopf witness \cite[Theorem~3.1]{LongHopf26}, with a change of circular normalization and a scalar factor. Let the bilinears be as in \eqref{eq:bilinears}, and put
\begin{equation}\label{eq:quartic}
 F_4=(R+X)\left(R^2Y-\frac{(2R+X)T^2}{2}\right),\qquad Q=X.
\end{equation}
Then, for every \(m\geq1\),
\begin{equation}\label{eq:quartic-identity}
 \E(F_4^m)=0,\qquad
 \E(QF_4^m)=\frac{(4m+2)!\,m!}{(2m+1)!!}\ne0.
\end{equation}
Here is a coefficient proof. On \(q=1\),
\[
 p_4=(1+x)\left(y-\frac{2+x}{2}t^2\right)
     =\frac{1+x}{2x}\bigl(1-t^2(1+x)^2\bigr).
\]
Define
\[
 J_m(U)=\int_0^U(1-v^2)^m\dd v,\qquad
 H_m(U)=U^m\int_0^1(1-s^2U^2)^m\dd s=U^{m-1}J_m(U).
\]
By \eqref{eq:ct},
\[
 L(p_4^m)=2^{-m}[u^m]H_m(1+u),\qquad
 L(xp_4^m)=2^{-m}[u^{m-1}]H_m(1+u).
\]
Since \(J_m'(U)=(1-U^2)^m\),
\(J_m(1+u)=J_m(1)+O(u^{m+1})\).
The first displayed coefficient is zero; the second is
\(2^{-m}J_m(1)=m!/(2m+1)!!\). Lemma~\ref{lem:bridge} proves
\eqref{eq:quartic-identity}, over every characteristic-zero field.
This example concerns a single fixed balanced bidegree four; it proves
exactly the minimum-degree assertion in Theorem~\ref{thm:image}.

\paragraph{Verification.}
All finite eliminations used in the proof are displayed in the article.
The accompanying short exact-arithmetic program regenerates the moment
identities, checks the two small polynomial obstructions, and verifies
the boundary-chart equations. It is a reproducibility aid; the geometric
reduction and all-order arguments are given here.

\paragraph{AI assistance and review.}
ChatGPT (OpenAI) assisted with mathematical arguments, literature checks,
exposition, proof review and verification code.

\begin{thebibliography}{9}
\bibitem{EWZ11}
A.~van den Essen, D.~Wright and W.~Zhao,
\emph{On the Image Conjecture}, J.~Algebra \textbf{340} (2011), 211--224.
\href{https://doi.org/10.1016/j.jalgebra.2011.04.036}{doi:10.1016/j.jalgebra.2011.04.036}.
\bibitem{Long26}
C.~D.~Long,
\emph{Counterexamples to the \(xz\)-Conjecture and the Mathieu Conjecture
for \(SU(2)\)}, arXiv:2607.19012v1 (2026).
\href{https://arxiv.org/abs/2607.19012v1}{arXiv:2607.19012v1}.
\bibitem{LongHopf26}
C.~D.~Long, \emph{The Mathieu Property for Compact Connected Lie Groups},
arXiv:2609.16178v1 (2026).
\href{https://arxiv.org/abs/2609.16178v1}{arXiv:2609.16178v1}.
\bibitem{Kas76}
M.~Kashiwara, \emph{B-functions and holonomic systems: rationality of
roots of B-functions}, Invent.~Math. \textbf{38} (1976), 33--53.
\href{https://www.kurims.kyoto-u.ac.jp/~kenkyubu/kashiwara/Bfunct.pdf}{Author's copy}.
\bibitem{CDM24}
Q.~Chen, B.~Dirks and M.~Musta\c t\u a,
\emph{An introduction to \(V\)-filtrations}, arXiv:2402.07816v1 (2024).
\href{https://arxiv.org/abs/2402.07816v1}{arXiv:2402.07816v1}.
\end{thebibliography}
\end{document}